% answers.tex —— 2013 年·湖北卷·文科 · 参考答案与详细解析片段
% 由 tools/generate_projects.py 依来源总账自 2013/2013_数学真题_湖北文_试题.md + 2013/2013_数学真题_湖北文_解析正文.md + 2013/2013_数学真题_湖北文_答案要点.md 生成，勿手改（改源 MD 后重跑）

\begin{answer}
\gksolhead{1}{}

\textbf{题目}：

已知全集 $U=\{1,2,3,4,5\}$，集合 $A=\{1,2\}$，$B=\{2,3,4\}$，则 $B\cap\complement_U A=$\mbox{（　　）}

（A）$\{2\}$；（B）$\{3,4\}$；（C）$\{1,4,5\}$；（D）$\{2,3,4,5\}$。

\Needspace{4\baselineskip}
\textbf{解析}
\Needspace{5\baselineskip}
由 $U=\{1,2,3,4,5\}$，$A=\{1,2\}$ 得
\nopagebreak[4]
\[\complement_U A=\{3,4,5\},\]
\Needspace{5\baselineskip}
故
\nopagebreak[4]
\[B\cap\complement_U A=\{2,3,4\}\cap\{3,4,5\}=\{3,4\}.\]
选 B。

\gkanswer{答案：B。}

\end{answer}
\begin{answer}
\gksolhead{2}{}

\textbf{题目}：

已知 $0<\theta<\dfrac\pi4$，则双曲线 $C_1:\dfrac{x^2}{\sin^2\theta}-\dfrac{y^2}{\cos^2\theta}=1$ 与 $C_2:\dfrac{y^2}{\cos^2\theta}-\dfrac{x^2}{\sin^2\theta}=1$ 的\mbox{（　　）}

（A）实轴长相等；（B）虚轴长相等；（C）离心率相等；（D）焦距相等。

\Needspace{4\baselineskip}
\textbf{解析}
\Needspace{5\baselineskip}
对 $C_1$，实半轴 $a_1=\sin\theta$，虚半轴 $b_1=\cos\theta$，故
\nopagebreak[4]
\[c_1^2=a_1^2+b_1^2=\sin^2\theta+\cos^2\theta=1.\]
\Needspace{5\baselineskip}
对 $C_2$，实半轴 $a_2=\cos\theta$，虚半轴 $b_2=\sin\theta$，故
\nopagebreak[4]
\[c_2^2=a_2^2+b_2^2=\cos^2\theta+\sin^2\theta=1.\]
因此 $c_1=c_2=1$，两双曲线的焦距都为 $2$。又 $0<\theta<\dfrac\pi4$ 时 $\sin\theta\neq\cos\theta$，实轴长、虚轴长、离心率均一般不相等。选 D。

\gkanswer{答案：D。}

\end{answer}
\begin{answer}
\gksolhead{3}{}

\textbf{题目}：

在一次跳伞训练中，甲、乙两位学员各跳一次。设命题 $p$ 是“甲降落在指定范围”，$q$ 是“乙降落在指定范围”，则命题“至少有一位学员没有降落在指定范围”可表示为\mbox{（　　）}

（A）$(\neg p)\lor(\neg q)$；（B）$p\lor(\neg q)$；（C）$(\neg p)\land(\neg q)$；（D）$p\lor q$。

\Needspace{4\baselineskip}
\textbf{解析}
\Needspace{5\baselineskip}
“至少有一位学员没有降落在指定范围”即“甲没有降落或乙没有降落”，也就是 $\neg p$ 与 $\neg q$ 至少一个成立：
\nopagebreak[4]
\[(\neg p)\lor(\neg q).\]
选 A。

\gkanswer{答案：A。}

\end{answer}
\begin{answer}
\gksolhead{4}{}

\textbf{题目}：

四名同学根据各自的样本数据研究变量 $x$，$y$ 之间的相关关系，并求得回归直线方程，分别得到以下四个结论：

① $y$ 与 $x$ 负相关且 $\hat y=2.347x-6.423$；

② $y$ 与 $x$ 负相关且 $\hat y=-3.476x+5.648$；

③ $y$ 与 $x$ 正相关且 $\hat y=5.437x+8.493$；

④ $y$ 与 $x$ 正相关且 $\hat y=-4.326x-4.578$。

其中一定不正确的结论的序号是\mbox{（　　）}

（A）①②；（B）②③；（C）③④；（D）①④。

\Needspace{4\baselineskip}
\textbf{解析}
正相关时回归直线的斜率应为正，负相关时斜率应为负。逐条判断：

\begin{itemize}
\item ① 说负相关，但斜率 $2.347>0$，矛盾，一定不正确；
\item ② 负相关且斜率 $-3.476<0$，可以成立；
\item ③ 正相关且斜率 $5.437>0$，可以成立；
\item ④ 说正相关，但斜率 $-4.326<0$，矛盾，一定不正确。
\end{itemize}

故一定不正确的是①④。选 D。

\gkanswer{答案：D。}

\end{answer}
\begin{answer}
\gksolhead{5}{}

\textbf{题目}：

小明骑车上学，开始时匀速行驶，途中因交通堵塞停留了一段时间后，为了赶时间加快速度行驶。与以上事件吻合得最好的图象是\mbox{（　　）}
\nopagebreak[4]
\begin{center}% 2013 湖北卷文科第 5 题：四个“距学校的距离—时间”图象。
\begin{tikzpicture}[font=\footnotesize, >=Stealth, line width=0.5pt]
  % (A)：距离从 0 开始增大，先匀速，后短暂平缓，再加速上升
  \begin{scope}[shift={(0,0)}]
    \draw[->] (0,0) -- (0,2.05) node[above] {距学校的距离};
    \draw[->] (0,0) -- (3.6,0) node[below right] {时间};
    \node[below left] at (0,0) {$O$};
    \draw[thick] (0,0) -- (1.25,0.58) -- (2.05,0.66)
      .. controls (2.45,0.68) and (2.70,0.95) .. (2.98,1.58);
    \node at (1.8,-0.55) {（A）};
  \end{scope}
  % (B)：先较快下降，短暂变缓，再匀速下降
  \begin{scope}[shift={(5.75,0)}]
    \draw[->] (0,0) -- (0,2.05) node[above] {距学校的距离};
    \draw[->] (0,0) -- (3.6,0) node[below right] {时间};
    \node[below left] at (0,0) {$O$};
    \draw[thick] (0,1.58) -- (1.55,1.08) -- (2.10,1.02) -- (3.25,0.16);
    \node at (1.8,-0.55) {（B）};
  \end{scope}
  % (C)：匀速下降，中途停留（水平），再加速下降
  \begin{scope}[shift={(0,-3.75)}]
    \draw[->] (0,0) -- (0,2.05) node[above] {距学校的距离};
    \draw[->] (0,0) -- (3.6,0) node[below right] {时间};
    \node[below left] at (0,0) {$O$};
    \draw[thick] (0,1.60) -- (1.20,1.10) -- (2.05,1.08)
      .. controls (2.42,1.06) and (2.62,0.72) .. (2.95,0.05);
    \node at (1.8,-0.55) {（C）};
  \end{scope}
  % (D)：匀速下降，随后直接加速下降（无停留）
  \begin{scope}[shift={(5.75,-3.75)}]
    \draw[->] (0,0) -- (0,2.05) node[above] {距学校的距离};
    \draw[->] (0,0) -- (3.6,0) node[below right] {时间};
    \node[below left] at (0,0) {$O$};
    \draw[thick] (0,1.60) -- (1.85,0.92) -- (2.85,0.05);
    \node at (1.8,-0.55) {（D）};
  \end{scope}
\end{tikzpicture}
\end{center}
\nopagebreak[4]
\Needspace{4\baselineskip}
\textbf{解析}
“距学校的距离”随出发后时间的增加而减小，故图象必须整体下降，排除（A）。

\begin{itemize}
\item 开始匀速行驶：第一段应为直线段，斜率为负且不变；
\item 交通堵塞停留：距离不随时间变化，图象应有一段水平线段；
\item 停留后加快速度：之后下降更快，直线段更陡（或曲线向下弯曲更急）。
\end{itemize}

（B）只是中途变缓，没有停留；（D）由匀速直接变为加速，没有停留；（C）先匀速下降，再水平停留，最后加速下降，与事件吻合。选 C。

\gkanswer{答案：C。}

\end{answer}
\begin{answer}
\gksolhead{6}{}

\textbf{题目}：

将函数 $y=\sqrt3\cos x+\sin x$（$x\in\mathbf R$）的图象向左平移 $m$（$m>0$）个单位长度后，所得到的图象关于 $y$ 轴对称，则 $m$ 的最小值是\mbox{（　　）}

（A）$\dfrac\pi{12}$；（B）$\dfrac\pi6$；（C）$\dfrac\pi3$；（D）$\dfrac{5\pi}6$。

\Needspace{4\baselineskip}
\textbf{解析}
\nopagebreak[4]
\[y=\sqrt3\cos x+\sin x=2\left(\frac{\sqrt3}2\cos x+\frac12\sin x\right)=2\cos\left(x-\frac\pi6\right).\]
\Needspace{5\baselineskip}
向左平移 $m$ 个单位后得
\nopagebreak[4]
\[y=2\cos\left(x+m-\frac\pi6\right).\]
\Needspace{5\baselineskip}
其图象关于 $y$ 轴对称，即函数为偶函数，须
\nopagebreak[4]
\[m-\frac\pi6=k\pi\quad(k\in\mathbf Z),\]
故 $m=\dfrac\pi6+k\pi$。由 $m>0$，最小值为 $\dfrac\pi6$（$k=0$）。选 B。

\gkanswer{答案：B。}

\end{answer}
\begin{answer}
\gksolhead{7}{}

\textbf{题目}：

已知点 $A(-1,1)$，$B(1,2)$，$C(-2,-1)$，$D(3,4)$，则向量 $\overrightarrow{AB}$ 在 $\overrightarrow{CD}$ 方向上的投影为\mbox{（　　）}

（A）$\dfrac{3\sqrt2}2$；（B）$\dfrac{3\sqrt{15}}2$；（C）$-\dfrac{3\sqrt2}2$；（D）$-\dfrac{3\sqrt{15}}2$。

\Needspace{4\baselineskip}
\textbf{解析}
\nopagebreak[4]
\[\overrightarrow{AB}=(1-(-1),\,2-1)=(2,1),\qquad
\overrightarrow{CD}=(3-(-2),\,4-(-1))=(5,5).\]
\Needspace{5\baselineskip}
$\overrightarrow{AB}$ 在 $\overrightarrow{CD}$ 方向上的投影为
\nopagebreak[4]
\[\frac{\overrightarrow{AB}\cdot\overrightarrow{CD}}{|\overrightarrow{CD}|}
=\frac{2\times5+1\times5}{\sqrt{5^2+5^2}}
=\frac{15}{5\sqrt2}=\frac{3\sqrt2}2.\]
选 A。

\gkanswer{答案：A。}

\end{answer}
\begin{answer}
\gksolhead{8}{}

\textbf{题目}：

$x$ 为实数，$[x]$ 表示不超过 $x$ 的最大整数，则函数 $f(x)=x-[x]$ 在 $\mathbf R$ 上为\mbox{（　　）}

（A）奇函数；（B）偶函数；（C）增函数；（D）周期函数。

\Needspace{4\baselineskip}
\textbf{解析}
\Needspace{5\baselineskip}
设 $f(x)=x-[x]$，则 $0\leqslant f(x)<1$，且
\nopagebreak[4]
\[f(x+1)=(x+1)-[x+1]=x+1-([x]+1)=x-[x]=f(x),\]
故 $f$ 是以 $1$ 为周期的周期函数。

又 $f(0)=0$，$f\!\left(\dfrac12\right)=\dfrac12$，$f\!\left(-\dfrac12\right)=\dfrac12$，故 $f$ 不是奇函数；$f\!\left(\dfrac13\right)=\dfrac13$，$f\!\left(-\dfrac13\right)=\dfrac23$，故 $f$ 不是偶函数；$f(0.9)=0.9>f(1)=0$，故 $f$ 不是增函数。选 D。

\gkanswer{答案：D。}

\end{answer}
\begin{answer}
\gksolhead{9}{}

\textbf{题目}：

某旅行社租用 $A$，$B$ 两种型号的客车安排 $900$ 名客人旅行，$A$，$B$ 两种车辆的载客量分别为 $36$ 人和 $60$ 人，租金分别为 $1600$ 元/辆和 $2400$ 元/辆，旅行社要求租车总数不超过 $21$ 辆，且 $B$ 型车不多于 $A$ 型车 $7$ 辆，则租金最少为\mbox{（　　）}

（A）$31200$ 元；（B）$36000$ 元；（C）$36800$ 元；（D）$38400$ 元。

\Needspace{4\baselineskip}
\textbf{解析}
\Needspace{5\baselineskip}
设租用 $A$ 型车 $x$ 辆，$B$ 型车 $y$ 辆（$x,y\in\mathbf N$），则
\nopagebreak[4]
\[\begin{cases}
36x+60y\geqslant900,\\
x+y\leqslant21,\\
y\leqslant x+7.
\end{cases}\]
\Needspace{5\baselineskip}
即
\nopagebreak[4]
\[\begin{cases}
3x+5y\geqslant75,\\
x+y\leqslant21,\\
y-x\leqslant7.
\end{cases}\]
租金为 $1600x+2400y=800(2x+3y)$。

可行域的顶点为直线 $3x+5y=75$ 与 $y=x+7$ 的交点 $(5,12)$、直线 $3x+5y=75$ 与 $x+y=21$ 的交点 $(15,6)$、直线 $y=x+7$ 与 $x+y=21$ 的交点 $(7,14)$。分别计算：

\begin{itemize}
\item $(5,12)$：租金 $5\times1600+12\times2400=36800$，载客 $36\times5+60\times12=900$；
\item $(15,6)$：租金 $15\times1600+6\times2400=38400$；
\item $(7,14)$：租金 $7\times1600+14\times2400=44800$。
\end{itemize}

故租金最少为 $36800$ 元，此时租 $A$ 型车 $5$ 辆、$B$ 型车 $12$ 辆。选 C。

\gkanswer{答案：C。}

\end{answer}
\begin{answer}
\gksolhead{10}{}

\textbf{题目}：

已知函数 $f(x)=x(\ln x-ax)$ 有两个极值点，则实数 $a$ 的取值范围是\mbox{（　　）}

（A）$(-\infty,0)$；（B）$\left(0,\dfrac12\right)$；（C）$(0,1)$；（D）$(0,+\infty)$。

\Needspace{4\baselineskip}
\textbf{解析}
\Needspace{5\baselineskip}
$f(x)=x\ln x-ax^2$ 的定义域为 $x>0$，且
\nopagebreak[4]
\[f'(x)=\ln x+1-2ax.\]
\Needspace{5\baselineskip}
令
\nopagebreak[4]
\[\varphi(x)=\ln x+1-2ax,\]
则 $f$ 有两个极值点当且仅当 $\varphi(x)=0$ 有两个不同的正根，且 $\varphi$ 在这些根两侧变号。

\Needspace{5\baselineskip}
求导得
\nopagebreak[4]
\[\varphi'(x)=\frac1x-2a.\]

\begin{itemize}
\item 若 $a\leqslant0$，则 $\varphi'(x)>0$，$\varphi$ 在 $(0,+\infty)$ 上严格递增，至多一个零点，不合题意。
\item 若 $a>0$，则 $\varphi$ 在 $\left(0,\dfrac1{2a}\right)$ 上递增，在 $\left(\dfrac1{2a},+\infty\right)$ 上递减，最大值为
\end{itemize}
\nopagebreak[4]
\[\varphi\left(\frac1{2a}\right)=\ln\frac1{2a}+1-2a\cdot\frac1{2a}=-\ln(2a).\]
\Needspace{5\baselineskip}
又 $x\to0^+$ 时 $\varphi(x)\to-\infty$，$x\to+\infty$ 时 $\varphi(x)\to-\infty$，故 $\varphi$ 有两个不同正根当且仅当
\nopagebreak[4]
\[-\ln(2a)>0\iff 2a<1\iff 0<a<\frac12.\]

当 $0<a<\dfrac12$ 时，两个零点把定义域分成三段，$f'$ 的符号依次为负、正、负，故 $f$ 先减、再增、再减，确有两个极值点。因此 $a\in\left(0,\dfrac12\right)$。选 B。

\gkanswer{答案：B。}

\end{answer}
\begin{answer}
\gksolhead{11}{}

\textbf{题目}：

$\mathrm i$ 为虚数单位，设复数 $z_1$，$z_2$ 在复平面内对应的点关于原点对称，若 $z_1=2-3\mathrm i$，则 $z_2=$\gkblank{}。

\Needspace{4\baselineskip}
\textbf{解析}
\Needspace{5\baselineskip}
复数 $z_1=2-3\mathrm i$ 对应点 $(2,-3)$。关于原点对称的点为 $(-2,3)$，故
\nopagebreak[4]
\[z_2=-2+3\mathrm i.\]

\gkanswer{答案：$-2+3\mathrm i$。}

\end{answer}
\begin{answer}
\gksolhead{12}{}

\textbf{题目}：

\Needspace{5\baselineskip}
某学员在一次射击测试中射靶 $10$ 次，命中环数如下：

\nopagebreak[4]
\[7,\ 8,\ 7,\ 9,\ 5,\ 4,\ 9,\ 10,\ 7,\ 4\]

则：

（1） 平均命中环数为\gkblank{}；

（2） 命中环数的标准差为\gkblank{}。

\Needspace{4\baselineskip}
\textbf{解析}
\Needspace{5\baselineskip}
\textbf{（1）} 命中环数的总和为
\nopagebreak[4]
\[7+8+7+9+5+4+9+10+7+4=70,\]
\Needspace{5\baselineskip}
故平均命中环数为
\nopagebreak[4]
\[\bar x=\frac{70}{10}=7.\]

\Needspace{5\baselineskip}
\textbf{（2）} 各数据与平均数的偏差为
\nopagebreak[4]
\[0,\ 1,\ 0,\ 2,\ -2,\ -3,\ 2,\ 3,\ 0,\ -3,\]
\Needspace{5\baselineskip}
偏差平方和为
\nopagebreak[4]
\[0^2+1^2+0^2+2^2+(-2)^2+(-3)^2+2^2+3^2+0^2+(-3)^2=40.\]
\Needspace{5\baselineskip}
故方差为 $\dfrac{40}{10}=4$，标准差为
\nopagebreak[4]
\[s=\sqrt4=2.\]

\gkanswer{答案（1）：$7$。}

\gkanswer{答案（2）：$2$。}

\end{answer}
\begin{answer}
\gksolhead{13}{}

\textbf{题目}：

阅读如图所示的程序框图，运行相应的程序。若输入 $m$ 的值为 $2$，则输出的结果 $i=$\gkblank{}。
\nopagebreak[4]
\begin{center}% 2013 湖北卷文科第 13 题：程序框图。
\begin{tikzpicture}[font=\footnotesize, >=Stealth, line width=0.45pt,
  node distance=4.2mm,
  box/.style={draw, minimum height=1.7em, inner xsep=5pt},
  term/.style={draw, rounded corners=5pt, minimum height=1.7em, inner xsep=8pt},
  io/.style={draw, trapezium, trapezium left angle=70, trapezium right angle=110,
             minimum height=1.7em, inner xsep=5pt},
  dec/.style={draw, diamond, aspect=2.1, inner sep=1pt}]
  \node[term] (start) {开始};
  \node[io, below=5mm of start] (input) {输入 $m$};
  \node[box, below=5mm of input] (init) {$A=1$，$B=1$，$i=0$};
  \node[box, below=6mm of init] (inc) {$i=i+1$};
  \node[box, below=4.5mm of inc] (am) {$A=A\times m$};
  \node[box, below=4.5mm of am] (bi) {$B=B\times i$};
  \node[dec, below=6mm of bi] (cond) {$A<B?$};
  \node[io, below=7mm of cond] (output) {输出 $i$};
  \node[term, below=5mm of output] (end) {结束};
  \draw[->] (start) -- (input);
  \draw[->] (input) -- (init);
  \draw[->] (init) -- (inc);
  \draw[->] (inc) -- (am);
  \draw[->] (am) -- (bi);
  \draw[->] (bi) -- (cond);
  \draw[->] (cond) -- node[right] {是} (output);
  \draw[->] (output) -- (end);
  % “否”分支：从判断框右侧引出，向上回到 i=i+1 之前的入口
  \draw[->] (cond.east) -- node[above] {否} ++(1.15,0)
    |- ($(inc.north)!0.5!(init.south)$);
\end{tikzpicture}
\end{center}
\nopagebreak[4]
\Needspace{4\baselineskip}
\textbf{解析}
输入 $m=2$，初值 $A=1$，$B=1$，$i=0$。每轮先执行 $i=i+1$，再执行 $A=A\times m$，$B=B\times i$，最后判断 $A<B$：

\begin{center}
\begin{tabular}{lllll}
\toprule
轮次 & $i$ & $A$ & $B$ & $A<B$ \\
\midrule
1 & 1 & 2 & 1 & 否 \\
2 & 2 & 4 & 2 & 否 \\
3 & 3 & 8 & 6 & 否 \\
4 & 4 & 16 & 24 & 是 \\
\bottomrule
\end{tabular}
\end{center}

第 4 轮判断为“是”，输出 $i=4$。

\gkanswer{答案：$4$。}

\end{answer}
\begin{answer}
\gksolhead{14}{}

\textbf{题目}：

已知圆 $O:x^2+y^2=5$，直线 $l:x\cos\theta+y\sin\theta=1\left(0<\theta<\dfrac\pi2\right)$。设圆 $O$ 上到直线 $l$ 的距离等于 $1$ 的点的个数为 $k$，则 $k=$\gkblank{}。

\Needspace{4\baselineskip}
\textbf{解析}
\Needspace{5\baselineskip}
原点 $O$ 到直线 $l:x\cos\theta+y\sin\theta=1$ 的距离为
\nopagebreak[4]
\[d=\frac{|0+0-1|}{\sqrt{\cos^2\theta+\sin^2\theta}}=1.\]
\Needspace{5\baselineskip}
到直线 $l$ 的距离等于 $1$ 的点，位于与 $l$ 平行且到 $l$ 距离为 $1$ 的两条直线上。这两条直线到圆心 $O$ 的距离分别为
\nopagebreak[4]
\[|1-1|=0,\qquad 1+1=2.\]
\Needspace{5\baselineskip}
其中过圆心的一条与圆 $x^2+y^2=5$ 交于 $2$ 个点；另一条到圆心的距离 $2<\sqrt5$，也与圆交于 $2$ 个点。因此
\nopagebreak[4]
\[k=2+2=4.\]

\gkanswer{答案：$4$。}

\end{answer}
\begin{answer}
\gksolhead{15}{}

\textbf{题目}：

在区间 $[-2,4]$ 上随机地取一个数 $x$，若 $x$ 满足 $|x|\leqslant m$ 的概率为 $\dfrac56$，则 $m=$\gkblank{}。

\Needspace{4\baselineskip}
\textbf{解析}
\Needspace{5\baselineskip}
区间 $[-2,4]$ 的长度为 $6$。事件 $|x|\leqslant m$ 对应区间 $[-m,m]$ 与 $[-2,4]$ 的交集，由几何概型，
\nopagebreak[4]
\[\frac{\text{\text{交集长度}}}6=\frac56,\]
故交集长度为 $5$。

\Needspace{5\baselineskip}
若 $0\leqslant m\leqslant2$，交集为 $[-m,m]$，长度为 $2m=5$，得 $m=2.5>2$，不合；若 $2\leqslant m\leqslant4$，交集为 $[-2,m]$，长度为 $m+2=5$，得 $m=3$，符合。故
\nopagebreak[4]
\[m=3.\]

\gkanswer{答案：$3$。}

\end{answer}
\begin{answer}
\gksolhead{16}{}

\textbf{题目}：

我国古代数学名著《数书九章》中有“天池盆测雨”题：在下雨时，用一个圆台形的天池盆接雨水，天池盆盆口直径为二尺八寸，盆底直径为一尺二寸，盆深一尺八寸，若盆中积水深九寸，则平地降雨量是\gkblank{}寸。

（注：① 平地降雨量等于盆中积水体积除以盆口面积；② 一尺等于十寸）

\Needspace{4\baselineskip}
\textbf{解析}
\Needspace{5\baselineskip}
一尺等于十寸。盆口半径 $R=1.4$ 尺，盆底半径 $r=0.6$ 尺，盆深 $1.8$ 尺。积水深 $0.9$ 尺，恰为盆深的一半，故水面半径为
\nopagebreak[4]
\[r'=0.6+\frac{0.9}{1.8}(1.4-0.6)=0.6+0.4=1.0\ \text{\text{尺}}.\]
\Needspace{5\baselineskip}
积水体积为圆台体积
\nopagebreak[4]
\[V_{\text{\text{水}}}=\frac13\pi\cdot0.9\cdot\left(1.0^2+1.0\times0.6+0.6^2\right)
=0.3\pi\times1.96=0.588\pi\ \text{\text{立方尺}}.\]
\Needspace{5\baselineskip}
盆口面积为
\nopagebreak[4]
\[S_{\text{\text{口}}}=\pi R^2=\pi\times1.4^2=1.96\pi\ \text{\text{平方尺}}.\]
\Needspace{5\baselineskip}
故平地降雨量为
\nopagebreak[4]
\[\frac{V_{\text{\text{水}}}}{S_{\text{\text{口}}}}=\frac{0.588\pi}{1.96\pi}=0.3\ \text{\text{尺}}=3\ \text{\text{寸}}.\]

\gkanswer{答案：$3$。}

\end{answer}
\begin{answer}
\gksolhead{17}{}

\textbf{题目}：

在平面直角坐标系中，若点 $P(x,y)$ 的坐标 $x$，$y$ 均为整数，则称点 $P$ 为格点，若一个多边形的顶点全是格点，则称该多边形为格点多边形，格点多边形的面积为 $S$，其内部的格点数记为 $N$，边界上的格点数记为 $L$。例如图中 $\triangle ABC$ 是格点三角形，对应的 $S=1$，$N=0$，$L=4$。
\nopagebreak[4]
\begin{center}% 2013 湖北卷文科第 17 题：格点四边形 DEFG 与格点三角形 ABC。
\begin{tikzpicture}[font=\footnotesize, >=Stealth, line width=0.45pt, scale=0.9]
  % 网格
  \draw[gray!45, very thin] (0,0) grid (6,6);
  % 坐标轴
  \draw[->] (0,0) -- (6.55,0) node[below right] {$x$};
  \draw[->] (0,0) -- (0,6.55) node[left] {$y$};
  \node[below left] at (0,0) {$O$};
  \foreach \x in {1,...,5} {
    \draw (\x,0) -- (\x,-0.06);
    \node[below] at (\x,-0.09) {$\x$};
  }
  \foreach \y in {1,...,5} {
    \draw (0,\y) -- (-0.06,\y);
    \node[left] at (-0.09,\y) {$\y$};
  }
  % 格点三角形 ABC
  \draw[thick] (1,2) -- (1,1) -- (3,1) -- cycle;
  \fill (1,2) circle (1.6pt) (1,1) circle (1.6pt) (3,1) circle (1.6pt);
  \node[above left] at (1,2) {$B$};
  \node[below left] at (1,1) {$C$};
  \node[below] at (3,1) {$A$};
  % 格点四边形 DEFG
  \draw[thick] (2,3) -- (4,3) -- (2,5) -- (1,4) -- cycle;
  \fill (2,3) circle (1.6pt) (4,3) circle (1.6pt) (2,5) circle (1.6pt) (1,4) circle (1.6pt);
  \node[below left] at (2,3) {$D$};
  \node[right] at (4,3) {$E$};
  \node[above] at (2,5) {$F$};
  \node[left] at (1,4) {$G$};
  % 图中的其余格点
  \fill (2,4) circle (1.3pt) (3,4) circle (1.3pt) (3,3) circle (1.3pt);
\end{tikzpicture}
\end{center}
\nopagebreak[4]
（1） 图中格点四边形 $DEFG$ 对应的 $S$，$N$，$L$ 分别是\gkblank{}；

（2） 已知格点多边形的面积可表示为 $S=aN+bL+c$，其中 $a$，$b$，$c$ 为常数。若某格点多边形对应的 $N=71$，$L=18$，则 $S=$\gkblank{}。（用数值作答）

\Needspace{4\baselineskip}
\textbf{解析}
\Needspace{5\baselineskip}
\textbf{（1）} 由图，四边形 $DEFG$ 的顶点为
\nopagebreak[4]
\[D(2,3),\quad E(4,3),\quad F(2,5),\quad G(1,4).\]
\Needspace{5\baselineskip}
用鞋带公式求面积：
\nopagebreak[4]
\[\begin{aligned}
S&=\frac12\Bigl|2\cdot3+4\cdot5+2\cdot4+1\cdot3\\
&\qquad -3\cdot4-3\cdot2-5\cdot1-4\cdot2\Bigr|\\
&=\frac12\left|6+20+8+3-12-6-5-8\right|=\frac12\times6=3.
\end{aligned}\]

\Needspace{5\baselineskip}
边界格点数按每条边上的格点数计算：
\nopagebreak[4]
\[L=\gcd(2,0)+\gcd(2,2)+\gcd(1,1)+\gcd(1,1)=2+2+1+1=6.\]
\Needspace{5\baselineskip}
由 Pick 定理 $S=N+\dfrac L2-1$，得
\nopagebreak[4]
\[N=S-\frac L2+1=3-3+1=1.\]
故 $S=3$，$N=1$，$L=6$。

\Needspace{5\baselineskip}
\textbf{（2）} 题设 $S=aN+bL+c$ 对一切格点多边形成立，而 Pick 定理给出
\nopagebreak[4]
\[S=N+\frac12L-1.\]
\Needspace{5\baselineskip}
对照系数得
\nopagebreak[4]
\[a=1,\qquad b=\frac12,\qquad c=-1.\]
（用题给 $\triangle ABC$ 验证：$N=0$，$L=4$，$S=1$，确有 $0+\dfrac12\times4-1=1$；用上面的 $DEFG$ 验证：$1+\dfrac12\times6-1=3$。）

\Needspace{5\baselineskip}
当 $N=71$，$L=18$ 时，
\nopagebreak[4]
\[S=71+\frac12\times18-1=71+9-1=79.\]

\gkanswer{答案（1）：$S=3$，$N=1$，$L=6$。}

\gkanswer{答案（2）：$79$。}

\end{answer}
\begin{answer}
\gksolhead{18}{}

\textbf{题目}：

在 $\triangle ABC$ 中，角 $A$，$B$，$C$ 对应的边分别是 $a$，$b$，$c$，已知 $\cos2A-3\cos(B+C)=1$。

（1） 求角 $A$ 的大小；

（2） 若 $\triangle ABC$ 的面积 $S=5\sqrt3$，$b=5$，求 $\sin B\sin C$ 的值。

\Needspace{4\baselineskip}
\textbf{解析}
\Needspace{5\baselineskip}
\textbf{（1）} 因为 $A+B+C=\pi$，所以
\nopagebreak[4]
\[\cos(B+C)=\cos(\pi-A)=-\cos A.\]
\Needspace{5\baselineskip}
代入已知式得
\nopagebreak[4]
\[\cos2A+3\cos A=1.\]
\Needspace{5\baselineskip}
由 $\cos2A=2\cos^2A-1$，
\nopagebreak[4]
\[2\cos^2A-1+3\cos A=1,\]
\Needspace{5\baselineskip}
即
\nopagebreak[4]
\[2\cos^2A+3\cos A-2=0,\]
\Needspace{5\baselineskip}
分解得
\nopagebreak[4]
\[(2\cos A-1)(\cos A+2)=0.\]
\Needspace{5\baselineskip}
因 $|\cos A|\leqslant1$，故 $\cos A=-2$ 不可能，只有
\nopagebreak[4]
\[\cos A=\frac12,\]
\Needspace{5\baselineskip}
从而
\nopagebreak[4]
\[A=\frac\pi3.\]

\Needspace{5\baselineskip}
\textbf{（2）} 三角形面积
\nopagebreak[4]
\[S=\frac12bc\sin A=\frac12\times5\times c\times\frac{\sqrt3}2=\frac{5\sqrt3}{4}c.\]
\Needspace{5\baselineskip}
由 $S=5\sqrt3$ 得
\nopagebreak[4]
\[\frac{5\sqrt3}{4}c=5\sqrt3\implies c=4.\]
\Needspace{5\baselineskip}
由余弦定理，
\nopagebreak[4]
\[a^2=b^2+c^2-2bc\cos A=25+16-2\times5\times4\times\frac12=41-20=21,\]
故 $a=\sqrt{21}$。

\Needspace{5\baselineskip}
由正弦定理 $\dfrac b{\sin B}=\dfrac c{\sin C}=\dfrac a{\sin A}$，得
\nopagebreak[4]
\[\sin B=\frac{b\sin A}{a},\qquad \sin C=\frac{c\sin A}{a}.\]
\Needspace{5\baselineskip}
于是
\nopagebreak[4]
\[\sin B\sin C=\frac{bc\sin^2A}{a^2}
=\frac{5\times4\times\left(\frac{\sqrt3}2\right)^2}{21}
=\frac{20\times\frac34}{21}
=\frac{15}{21}=\frac57.\]

\gkanswer{答案（1）：$A=\dfrac\pi3$。}

\gkanswer{答案（2）：$\dfrac57$。}

\end{answer}
\begin{answer}
\gksolhead{19}{}

\textbf{题目}：

已知 $S_n$ 是等比数列 $\{a_n\}$ 的前 $n$ 项和，$S_4$，$S_2$，$S_3$ 成等差数列，且 $a_2+a_3+a_4=-18$。

（1） 求数列 $\{a_n\}$ 的通项公式；

（2） 是否存在正整数 $n$，使得 $S_n\geqslant2013$？若存在，求出符合条件的所有的 $n$ 的集合；若不存在，说明理由。

\Needspace{4\baselineskip}
\textbf{解析}
\Needspace{5\baselineskip}
\textbf{（1）} 设等比数列的首项为 $a_1$，公比为 $q$（$q\neq0$）。则
\nopagebreak[4]
\[S_2=a_1(1+q),\qquad S_3=a_1(1+q+q^2),\qquad S_4=a_1(1+q+q^2+q^3).\]
\Needspace{5\baselineskip}
由 $S_4$，$S_2$，$S_3$ 成等差数列，
\nopagebreak[4]
\[2S_2=S_4+S_3,\]
\Needspace{5\baselineskip}
即
\nopagebreak[4]
\[2a_1(1+q)=a_1(2+2q+2q^2+q^3).\]
\Needspace{5\baselineskip}
若 $a_1=0$，则 $a_2=a_3=a_4=0$，与 $a_2+a_3+a_4=-18$ 矛盾，故 $a_1\neq0$。两边除以 $a_1$：
\nopagebreak[4]
\[2+2q=2+2q+2q^2+q^3,\]
\Needspace{5\baselineskip}
即
\nopagebreak[4]
\[q^2(q+2)=0.\]
所以 $q=0$ 或 $q=-2$。$q=0$ 时 $a_2=a_3=a_4=0$，和不为 $-18$，舍去；故 $q=-2$。

\Needspace{5\baselineskip}
此时
\nopagebreak[4]
\[a_2+a_3+a_4=a_1q+a_1q^2+a_1q^3
=a_1\left(-2+4-8\right)=-6a_1=-18,\]
\Needspace{5\baselineskip}
得 $a_1=3$。因此
\nopagebreak[4]
\[a_n=3(-2)^{n-1}\quad(n\in\mathbf N^*).\]

\Needspace{5\baselineskip}
\textbf{（2）} 由（1），等比数列的前 $n$ 项和为
\nopagebreak[4]
\[S_n=\frac{3\left[1-(-2)^n\right]}{1-(-2)}=1-(-2)^n.\]
\Needspace{5\baselineskip}
由 $S_n\geqslant2013$ 得
\nopagebreak[4]
\[1-(-2)^n\geqslant2013,\]
\Needspace{5\baselineskip}
即
\nopagebreak[4]
\[(-2)^n\leqslant-2012.\]

\Needspace{5\baselineskip}
若 $n$ 为偶数，则 $(-2)^n>0$，不等式不可能成立；故 $n$ 必须为奇数，此时 $(-2)^n=-2^n$，不等式化为
\nopagebreak[4]
\[2^n\geqslant2012.\]
\Needspace{5\baselineskip}
又
\nopagebreak[4]
\[2^{10}=1024<2012<2048=2^{11},\]
\Needspace{5\baselineskip}
故奇数 $n$ 须满足 $n\geqslant11$。因此所有符合条件的 $n$ 的集合为
\nopagebreak[4]
\[\{n\mid n=2k+1,\ k\in\mathbf N,\ k\geqslant5\}=\{11,13,15,\ldots\}.\]

\gkanswer{答案（1）：$a_n=3\times(-2)^{n-1}$。}

\gkanswer{答案（2）：$\{n\mid n=2k+1,\ k\in\mathbf N,\ k\geqslant5\}=\{11,13,15,\ldots\}$。}

\end{answer}
\begin{answer}
\gksolhead{20}{}

\textbf{题目}：

如图，某地质队自水平地面 $A$，$B$，$C$ 三处垂直向地下钻探，自 $A$ 点向下钻到 $A_1$ 处发现矿藏，再继续下钻到 $A_2$ 处后下面已无矿，从而得到在 $A$ 处正下方的矿层厚度为 $A_1A_2=d_1$，同样可得在 $B$，$C$ 处正下方的矿层厚度分别为 $B_1B_2=d_2$，$C_1C_2=d_3$，且 $d_1<d_2<d_3$，过 $AB$，$AC$ 的中点 $M$，$N$ 且与直线 $AA_2$ 平行的平面截多面体 $A_1B_1C_1-A_2B_2C_2$ 所得的截面 $DEFG$ 为该多面体的一个中截面，其面积记为 $S_{\text{\text{中}}}$。
\nopagebreak[4]
\begin{center}% 2013 湖北卷文科第 20 题：矿层多面体 A1B1C1-A2B2C2 与中截面 DEFG。
\begin{tikzpicture}[font=\footnotesize, >=Stealth, line width=0.5pt, scale=0.9]
  % 地面三角形 ABC（在水平面内）
  \coordinate (A) at (0,0);
  \coordinate (B) at (6,0);
  \coordinate (C) at (4.0,-1.30);
  % 矿层顶面三角形 A1B1C1
  \coordinate (A1) at (0,-1.70);
  \coordinate (B1) at (6,-1.90);
  \coordinate (C1) at (4.0,-3.00);
  % 矿层底面三角形 A2B2C2
  \coordinate (A2) at (0,-2.60);
  \coordinate (B2) at (6,-3.50);
  \coordinate (C2) at (4.0,-5.00);
  % 中点 M, N 与中截面顶点
  \coordinate (M) at (3,0);
  \coordinate (N) at (2,-0.65);
  \coordinate (D) at (3,-1.80);
  \coordinate (E) at (3,-3.05);
  \coordinate (G) at (2,-2.35);
  \coordinate (F) at (2,-3.80);

  % 地面与钻探竖线
  \draw (A) -- (B) -- (C) -- cycle;
  \draw (A) -- (A1) -- (A2);
  \draw (B) -- (B1) -- (B2);
  \draw (C) -- (C1) -- (C2);
  % 顶面、底面的边（背面用虚线）
  \draw[dashed] (A1) -- (B1);
  \draw (B1) -- (C1);
  \draw[dashed] (C1) -- (A1);
  \draw[dashed] (A2) -- (B2);
  \draw (B2) -- (C2);
  \draw[dashed] (C2) -- (A2);

  % 中截面 DEFG
  \fill[gray!18] (D) -- (E) -- (F) -- (G) -- cycle;
  \draw[very thick] (D) -- (E) -- (F) -- (G) -- cycle;
  % 截面所在平面与 M, N 的连线（辅助虚线）
  \draw[dashed] (M) -- (D);
  \draw[dashed] (N) -- (G);
  \draw[dashed] (M) -- (N);

  % 标注
  \node[above left] at (A) {$A$};
  \node[above right] at (B) {$B$};
  \node[right] at (C) {$C$};
  \node[left] at (A1) {$A_1$};
  \node[right] at (B1) {$B_1$};
  \node[right] at (C1) {$C_1$};
  \node[left] at (A2) {$A_2$};
  \node[right] at (B2) {$B_2$};
  \node[below] at (C2) {$C_2$};
  \node[above] at (M) {$M$};
  \node[left] at (N) {$N$};
  \node[right] at ($(D)+(0.30,-0.26)$) {$D$};
  \node[right] at ($(E)+(0.22,-0.10)$) {$E$};
  \node[left] at (F) {$F$};
  \node[left] at (G) {$G$};
\end{tikzpicture}
\end{center}
\nopagebreak[4]
（1） 证明：中截面 $DEFG$ 是梯形；

（2） 在 $\triangle ABC$ 中，记 $BC=a$，$BC$ 边上的高为 $h$，面积为 $S$。在估测三角形 $ABC$ 区域内正下方的矿藏储量（即多面体 $A_1B_1C_1-A_2B_2C_2$ 的体积 $V$）时，可用近似公式 $V_{\text{\text{估}}}=S_{\text{\text{中}}}\cdot h$ 来估算。已知 $V=\dfrac13(d_1+d_2+d_3)S$，试判断 $V_{\text{\text{估}}}$ 与 $V$ 的大小关系，并加以证明。

\Needspace{4\baselineskip}
\textbf{解析}
\Needspace{5\baselineskip}
\textbf{（1）} 因为截面所在平面与 $AA_2$ 平行，且该平面与平面 $ABB_1A_2$ 有一条公共点 $M$，所以它与平面 $ABB_1A_2$ 的交线过 $M$ 且平行于 $AA_2$。这条交线与 $A_1B_1$、$A_2B_2$ 分别交于 $D$、$E$，故
\nopagebreak[4]
\[DE\parallel AA_2.\]

\Needspace{5\baselineskip}
同理，该平面与平面 $ACC_1A_2$ 的交线过 $N$ 且平行于 $AA_2$，它与 $A_1C_1$、$A_2C_2$ 分别交于 $G$、$F$，故
\nopagebreak[4]
\[GF\parallel AA_2.\]

\Needspace{5\baselineskip}
于是
\nopagebreak[4]
\[DE\parallel GF.\]
四边形 $DEFG$ 有一组对边平行，所以它是梯形。

\Needspace{5\baselineskip}
\textbf{（2）} 由（1），$D$，$E$ 分别是 $A_1B_1$，$A_2B_2$ 的中点，$G$，$F$ 分别是 $A_1C_1$，$A_2C_2$ 的中点。在梯形 $A_1B_1B_2A_2$ 中，$DE$ 是中位线，且 $A_1A_2=d_1$，$B_1B_2=d_2$，故
\nopagebreak[4]
\[DE=\frac{d_1+d_2}{2}.\]
\Needspace{5\baselineskip}
同理，在梯形 $A_1C_1C_2A_2$ 中，
\nopagebreak[4]
\[GF=\frac{d_1+d_3}{2}.\]
\Needspace{5\baselineskip}
又 $DE$ 与 $GF$ 都是竖直方向，它们之间的距离等于 $M$ 与 $N$ 的距离。因 $M$，$N$ 分别为 $AB$，$AC$ 的中点，
\nopagebreak[4]
\[MN=\frac12BC=\frac a2.\]
\Needspace{5\baselineskip}
所以中截面面积
\nopagebreak[4]
\[S_{\text{\text{中}}}=\frac12(DE+GF)\cdot MN
=\frac12\left(\frac{d_1+d_2}{2}+\frac{d_1+d_3}{2}\right)\cdot\frac a2
=\frac{a(2d_1+d_2+d_3)}8.\]

\Needspace{5\baselineskip}
于是
\nopagebreak[4]
\[V_{\text{\text{估}}}=S_{\text{\text{中}}}\cdot h=\frac{ah(2d_1+d_2+d_3)}8.\]
\Needspace{5\baselineskip}
又 $S=\dfrac12ah$，故 $ah=2S$，从而
\nopagebreak[4]
\[V_{\text{\text{估}}}=\frac{2S(2d_1+d_2+d_3)}8=\frac{S(2d_1+d_2+d_3)}4.\]
\Needspace{5\baselineskip}
与
\nopagebreak[4]
\[V=\frac13(d_1+d_2+d_3)S\]
\Needspace{5\baselineskip}
作差：
\nopagebreak[4]
\[V-V_{\text{\text{估}}}
=S\left[\frac{d_1+d_2+d_3}3-\frac{2d_1+d_2+d_3}4\right]
=\frac{S(d_2+d_3-2d_1)}{12}.\]
\Needspace{5\baselineskip}
因为 $d_1<d_2<d_3$，所以
\nopagebreak[4]
\[d_2+d_3-2d_1>0,\]
\Needspace{5\baselineskip}
从而
\nopagebreak[4]
\[V-V_{\text{\text{估}}}>0,\]
\Needspace{5\baselineskip}
即
\nopagebreak[4]
\[V_{\text{\text{估}}}<V.\]

\gkanswer{答案（1）：证明见解析。}

\gkanswer{答案（2）：$V_{\text{\text{估}}}<V$。}

\end{answer}
\begin{answer}
\gksolhead{21}{}

\textbf{题目}：

设 $a>0$，$b>0$，已知函数 $f(x)=\dfrac{ax+b}{x+1}$。

（1） 当 $a\neq b$ 时，讨论函数 $f(x)$ 的单调性；

（2） 当 $x>0$ 时，称 $f(x)$ 为 $a$，$b$ 关于 $x$ 的加权平均数。

① 判断 $f(1)$，$f\left(\sqrt{\dfrac ba}\right)$，$f\left(\dfrac ba\right)$ 是否成等比数列，并证明 $f\left(\dfrac ba\right)\leqslant f\left(\sqrt{\dfrac ba}\right)$；

② $a$，$b$ 的几何平均数记为 $G$，称 $\dfrac{2ab}{a+b}$ 为 $a$，$b$ 的调和平均数，记为 $H$，若 $H\leqslant f(x)\leqslant G$，求 $x$ 的取值范围。

\Needspace{4\baselineskip}
\textbf{解析}
\Needspace{5\baselineskip}
\textbf{（1）} 函数的定义域为 $\{x\mid x\neq-1\}$。求导：
\nopagebreak[4]
\[f'(x)=\frac{a(x+1)-(ax+b)}{(x+1)^2}=\frac{a-b}{(x+1)^2}.\]
因为 $(x+1)^2>0$，所以 $f'(x)$ 的符号由 $a-b$ 决定。

当 $a>b$ 时，$f'(x)>0$，$f(x)$ 在 $(-\infty,-1)$ 和 $(-1,+\infty)$ 上分别单调递增；当 $a<b$ 时，$f'(x)<0$，$f(x)$ 在 $(-\infty,-1)$ 和 $(-1,+\infty)$ 上分别单调递减。

\Needspace{5\baselineskip}
\textbf{（2）①} 记 $t=\sqrt{\dfrac ba}$，则 $t>0$，$b=at^2$。直接计算：
\nopagebreak[4]
\[f(1)=\frac{a+b}{2},\]
\nopagebreak[4]
\[f\left(\sqrt{\frac ba}\right)=f(t)=\frac{at+b}{t+1}
=\frac{at+at^2}{t+1}=at=\sqrt{ab},\]
\nopagebreak[4]
\[f\left(\frac ba\right)=f(t^2)=\frac{at^2+b}{t^2+1}
=\frac{2at^2}{t^2+1}
=\frac{2ab}{a+b}.\]
\Needspace{5\baselineskip}
记
\nopagebreak[4]
\[A=\frac{a+b}{2},\qquad G=\sqrt{ab},\qquad H=\frac{2ab}{a+b}.\]
\Needspace{5\baselineskip}
则
\nopagebreak[4]
\[A\cdot H=\frac{a+b}{2}\cdot\frac{2ab}{a+b}=ab=G^2.\]
故 $f(1)$，$f\left(\sqrt{\dfrac ba}\right)$，$f\left(\dfrac ba\right)$ 成等比数列（按 $A$，$G$，$H$ 的顺序）。

\Needspace{5\baselineskip}
又由基本不等式
\nopagebreak[4]
\[a+b\geqslant2\sqrt{ab},\]
\Needspace{5\baselineskip}
得
\nopagebreak[4]
\[H=\frac{2ab}{a+b}\leqslant\frac{2ab}{2\sqrt{ab}}=\sqrt{ab}=G,\]
\Needspace{5\baselineskip}
即
\nopagebreak[4]
\[f\left(\frac ba\right)\leqslant f\left(\sqrt{\frac ba}\right).\]

\Needspace{5\baselineskip}
\textbf{②} 仍记 $t=\sqrt{\dfrac ba}$，$b=at^2$，$G=at$，$H=\dfrac{2at^2}{1+t^2}$。对 $x>0$，
\nopagebreak[4]
\[f(x)-H=\frac{ax+b}{x+1}-\frac{2at^2}{1+t^2}
=\frac{a(1-t^2)}{(1+t^2)(x+1)}(x-t^2),\]
\nopagebreak[4]
\[f(x)-G=\frac{ax+b}{x+1}-at
=\frac{a(1-t)(x-t)}{x+1}.\]

因为 $a>0$，$x+1>0$，所以

\begin{itemize}
\item 若 $b<a$，即 $0<t<1$，则 $H\leqslant f(x)$ 给出 $x\geqslant t^2$，$f(x)\leqslant G$ 给出 $x\leqslant t$，故
\end{itemize}
\nopagebreak[4]
\[\frac ba\leqslant x\leqslant\sqrt{\frac ba}.\]
\begin{itemize}
\item 若 $b>a$，即 $t>1$，则 $H\leqslant f(x)$ 给出 $x\leqslant t^2$，$f(x)\leqslant G$ 给出 $x\geqslant t$，故
\end{itemize}
\nopagebreak[4]
\[\sqrt{\frac ba}\leqslant x\leqslant\frac ba.\]
\begin{itemize}
\item 若 $a=b$，则 $t=1$，$G=H=a$，且 $f(x)=\dfrac{ax+a}{x+1}=a$，不等式对一切 $x>0$ 成立。
\end{itemize}

综上，当 $b<a$ 时 $x\in\left[\dfrac ba,\sqrt{\dfrac ba}\right]$；当 $b>a$ 时 $x\in\left[\sqrt{\dfrac ba},\dfrac ba\right]$；当 $a=b$ 时 $x\in(0,+\infty)$。

\gkanswer{答案（1）：当 $a>b$ 时，$f(x)$ 在 $(-\infty,-1)$ 和 $(-1,+\infty)$ 上单调递增；当 $a<b$ 时，$f(x)$ 在 $(-\infty,-1)$ 和 $(-1,+\infty)$ 上单调递减。}

\gkanswer{答案（2）：① 成等比数列，且 $f\left(\dfrac ba\right)\leqslant f\left(\sqrt{\dfrac ba}\right)$，证明见解析；② 当 $b<a$ 时，$\dfrac ba\leqslant x\leqslant\sqrt{\dfrac ba}$；当 $b>a$ 时，$\sqrt{\dfrac ba}\leqslant x\leqslant\dfrac ba$；当 $a=b$ 时，$x>0$。}

\end{answer}
\begin{answer}
\gksolhead{22}{}

\textbf{题目}：

如图，已知椭圆 $C_1$ 与 $C_2$ 的中心在坐标原点 $O$，长轴均为 $MN$ 且在 $x$ 轴上，短轴长分别为 $2m$，$2n$（$m>n$），过原点且不与 $x$ 轴重合的直线 $l$ 与 $C_1$，$C_2$ 的四个交点按纵坐标从大到小依次为 $A$，$B$，$C$，$D$，记 $\lambda=\dfrac mn$，$\triangle BDM$ 和 $\triangle ABN$ 的面积分别为 $S_1$ 和 $S_2$。
\nopagebreak[4]
\begin{center}% 2013 湖北卷文科第 22 题：共焦轴的两椭圆与过原点的直线 l。
\begin{tikzpicture}[font=\footnotesize, >=Stealth, line width=0.5pt, scale=1.0]
  \def\a{4}
  \def\m{2.4}
  \def\n{1.4}
  % 椭圆
  \draw[thick] (0,0) ellipse ({\a} and {\m});
  \draw[thick] (0,0) ellipse ({\a} and {\n});
  % 坐标轴
  \draw[->] (-4.8,0) -- (4.9,0) node[below right] {$x$};
  \draw[->] (0,-3.0) -- (0,3.1) node[left] {$y$};
  \node[below left] at (0,0) {$O$};
  % 四个交点（直线 y=0.75x）
  \coordinate (A) at (2.499,1.874);
  \coordinate (B) at (1.692,1.269);
  \coordinate (C) at (-1.692,-1.269);
  \coordinate (D) at (-2.499,-1.874);
  \coordinate (M) at (-\a,0);
  \coordinate (N) at (\a,0);
  % 直线 l
  \draw[thick] (-4.7,-3.525) -- (4.7,3.525);
  % 两个三角形 BDM 与 ABN 的边
  \draw (M) -- (B);
  \draw (B) -- (N);
  \draw (A) -- (N);
  \draw (M) -- (D);
  \draw (A) -- (B);
  % 交点与端点
  \fill (A) circle (1.6pt) (B) circle (1.6pt) (C) circle (1.6pt) (D) circle (1.6pt);
  \fill (M) circle (1.6pt) (N) circle (1.6pt);
  \node[fill=white, inner sep=1pt, above left] at (A) {$A$};
  \node[fill=white, inner sep=1pt, above left] at (B) {$B$};
  \node[fill=white, inner sep=1pt, above left] at (C) {$C$};
  \node[fill=white, inner sep=1pt, below right] at (D) {$D$};
  \node[fill=white, inner sep=1pt, above left] at (M) {$M$};
  \node[fill=white, inner sep=1pt, above right] at (N) {$N$};
\end{tikzpicture}
\end{center}
\nopagebreak[4]
（1） 当直线 $l$ 与 $y$ 轴重合时，若 $S_1=\lambda S_2$，求 $\lambda$ 的值；

（2） 当 $\lambda$ 变化时，是否存在与坐标轴不重合的直线 $l$，使得 $S_1=\lambda S_2$？并说明理由。

\Needspace{4\baselineskip}
\textbf{解析}
\Needspace{5\baselineskip}
设椭圆的半长轴为 $a$，则 $M(-a,0)$，$N(a,0)$。$C_1$，$C_2$ 的方程分别为
\nopagebreak[4]
\[C_1:\frac{x^2}{a^2}+\frac{y^2}{m^2}=1,\qquad
C_2:\frac{x^2}{a^2}+\frac{y^2}{n^2}=1.\]
直线 $l$ 过原点且不与 $x$ 轴重合，设其与 $x$ 轴正方向的夹角为 $\theta$，$s=\sin\theta>0$，$c=\cos\theta$。直线上的点可写成 $(tc,ts)$。

\Needspace{5\baselineskip}
对半短轴为 $\mu$（$\mu=m$ 或 $n$）的椭圆，交点满足
\nopagebreak[4]
\[t^2\left(\frac{c^2}{a^2}+\frac{s^2}{\mu^2}\right)=1,\]
\Needspace{5\baselineskip}
故交点到原点的距离为
\nopagebreak[4]
\[r_\mu=\frac1{\sqrt{\dfrac{c^2}{a^2}+\dfrac{s^2}{\mu^2}}}.\]
\Needspace{5\baselineskip}
因 $m>n$，故 $r_m>r_n$，且四个交点按纵坐标从大到小依次为
\nopagebreak[4]
\[A=(r_mc,r_ms),\quad B=(r_nc,r_ns),\quad C=(-r_nc,-r_ns),\quad D=(-r_mc,-r_ms).\]

\Needspace{5\baselineskip}
\textbf{（1）} 当 $l$ 与 $y$ 轴重合时，$c=0$，$s=1$，$r_m=m$，$r_n=n$，于是
\nopagebreak[4]
\[A(0,m),\quad B(0,n),\quad C(0,-n),\quad D(0,-m).\]
\Needspace{5\baselineskip}
此时 $BD=m+n$，$M$ 到 $y$ 轴的距离为 $a$，故
\nopagebreak[4]
\[S_1=S_{\triangle BDM}=\frac12(m+n)a.\]
\Needspace{5\baselineskip}
又 $AB=m-n$，$N$ 到 $y$ 轴的距离为 $a$，故
\nopagebreak[4]
\[S_2=S_{\triangle ABN}=\frac12(m-n)a.\]
\Needspace{5\baselineskip}
由 $S_1=\lambda S_2$ 及 $\lambda=\dfrac mn$，得
\nopagebreak[4]
\[\frac12(m+n)a=\frac mn\cdot\frac12(m-n)a,\]
\Needspace{5\baselineskip}
即
\nopagebreak[4]
\[n(m+n)=m(m-n).\]
\Needspace{5\baselineskip}
整理得
\nopagebreak[4]
\[m^2-2mn-n^2=0.\]
\Needspace{5\baselineskip}
两边除以 $n^2$，令 $\lambda=\dfrac mn>1$，得
\nopagebreak[4]
\[\lambda^2-2\lambda-1=0,\]
\Needspace{5\baselineskip}
解得
\nopagebreak[4]
\[\lambda=1\pm\sqrt2.\]
\Needspace{5\baselineskip}
因 $\lambda>1$，故
\nopagebreak[4]
\[\lambda=1+\sqrt2.\]

\Needspace{5\baselineskip}
\textbf{（2）} 对一般的直线 $l$，由上面的坐标，
\nopagebreak[4]
\[B-D=(r_m+r_n)(c,s),\qquad M-B=(-a-r_nc,-r_ns).\]
\Needspace{5\baselineskip}
用三角形面积的行列式公式，
\nopagebreak[4]
\[\begin{aligned}
S_1&=\frac12\left|(r_m+r_n)c\cdot(-r_ns)-(r_m+r_n)s\cdot(-a-r_nc)\right|\\
&=\frac12(r_m+r_n)a|s|.
\end{aligned}\]
\Needspace{5\baselineskip}
同理，
\nopagebreak[4]
\[\begin{aligned}
S_2&=\frac12\left|(r_n-r_m)c\cdot(-r_ms)-(r_n-r_m)s\cdot(a-r_mc)\right|\\
&=\frac12(r_m-r_n)a|s|.
\end{aligned}\]
\Needspace{5\baselineskip}
因此
\nopagebreak[4]
\[\frac{S_1}{S_2}=\frac{r_m+r_n}{r_m-r_n}.\]
\Needspace{5\baselineskip}
令
\nopagebreak[4]
\[\rho=\frac{r_m}{r_n}>1,\]
\Needspace{5\baselineskip}
则 $S_1=\lambda S_2$ 等价于
\nopagebreak[4]
\[\frac{\rho+1}{\rho-1}=\lambda,\]
\Needspace{5\baselineskip}
即
\nopagebreak[4]
\[\rho=\frac{\lambda+1}{\lambda-1}.\]

\Needspace{5\baselineskip}
另一方面，由 $r_\mu$ 的表达式，
\nopagebreak[4]
\[\rho^2=\frac{r_m^2}{r_n^2}
=\frac{\dfrac{c^2}{a^2}+\dfrac{s^2}{n^2}}{\dfrac{c^2}{a^2}+\dfrac{s^2}{m^2}}.\]
\Needspace{5\baselineskip}
当 $s^2$ 从 $0$ 增大到 $1$ 时，分子、分母都增大，但分子的增量系数 $\dfrac1{n^2}$ 大于分母的增量系数 $\dfrac1{m^2}$，故 $\rho^2$ 严格递增。于是
\nopagebreak[4]
\[\rho^2\in\left(1,\frac{m^2}{n^2}\right]=\left(1,\lambda^2\right],\]
\Needspace{5\baselineskip}
从而
\nopagebreak[4]
\[\rho\in(1,\lambda].\]
其中 $\rho=1$ 在 $s=0$（$l$ 与 $x$ 轴重合）时取到，$\rho=\lambda$ 在 $s=1$（$l$ 与 $y$ 轴重合）时取到。

\Needspace{5\baselineskip}
题目要求 $l$ 与坐标轴不重合，故 $0<s<1$，从而
\nopagebreak[4]
\[\rho\in(1,\lambda).\]
\Needspace{5\baselineskip}
因此存在这样的直线当且仅当
\nopagebreak[4]
\[1<\frac{\lambda+1}{\lambda-1}<\lambda.\]
\Needspace{5\baselineskip}
左半不等式显然成立。右半不等式等价于
\nopagebreak[4]
\[\lambda+1<\lambda(\lambda-1)=\lambda^2-\lambda,\]
\Needspace{5\baselineskip}
即
\nopagebreak[4]
\[\lambda^2-2\lambda-1>0.\]
\Needspace{5\baselineskip}
因 $\lambda>1$，解得
\nopagebreak[4]
\[\lambda>1+\sqrt2.\]

当 $\lambda>1+\sqrt2$ 时，由 $\rho$ 的连续性，存在 $\theta\in\left(0,\dfrac\pi2\right)$ 使 $\rho=\dfrac{\lambda+1}{\lambda-1}$（在 $\left(0,\dfrac\pi2\right)$ 内这样的 $\theta$ 唯一），从而存在满足条件的直线 $l$；由对称性，它关于 $x$ 轴的对称直线也满足条件。

当 $\lambda=1+\sqrt2$ 时，$\dfrac{\lambda+1}{\lambda-1}=\lambda$，只能由 $l$ 与 $y$ 轴重合取得，不符合“与坐标轴不重合”的要求；当 $1<\lambda<1+\sqrt2$ 时，$\dfrac{\lambda+1}{\lambda-1}>\lambda$，无解。

综上，当 $\lambda>1+\sqrt2$ 时存在；当 $1<\lambda\leqslant1+\sqrt2$ 时不存在。

\gkanswer{答案（1）：$\lambda=1+\sqrt2$。}

\gkanswer{答案（2）：当 $\lambda>1+\sqrt2$ 时存在；当 $1<\lambda\leqslant1+\sqrt2$ 时不存在。}

\end{answer}